\chapter{Introduction}

\section{Context: Natural Language Processing for Medical Documents}

The past two decades have witnessed a profound transformation in healthcare documentation, driven by the widespread adoption of Electronic Health Records (EHRs). These digital repositories now serve as the primary medium for capturing patient information, encompassing diagnoses, medications, treatment plans, laboratory results, imaging reports, and clinical narratives~\citep{shickel2018deep, wu2020deep}. The volume of healthcare data has grown at an unprecedented rate, with projections indicating that medical knowledge doubles approximately every 73 days~\citep{densen2011challenges}. This exponential growth presents both opportunities and challenges for healthcare systems seeking to leverage their data assets for improved patient care and clinical research.

A critical characteristic of EHR data is its heterogeneous nature, comprising both structured elements—such as coded diagnoses, vital signs, and laboratory values—and unstructured free-text narratives written by healthcare providers. Research consistently indicates that approximately 80\% of clinical data exists in unstructured format~\citep{kong2019managing, wu2022neural}, including physician notes, discharge summaries, radiology reports, and nursing assessments. This free-text component captures nuanced clinical information that structured data fields cannot adequately represent: the subtleties of symptom presentation, the reasoning behind clinical decisions, social determinants of health, and the evolution of patient conditions over time~\citep{wang2018clinical, dalianis2018clinical}.

The richness of clinical narratives makes them invaluable for understanding patient trajectories, yet their unstructured nature renders them largely inaccessible to traditional computational analysis. As~\citet{yang2022large} note, clinical notes provide a ``glimpse into the physician's brain,'' documenting not only what was observed but also the clinical reasoning process. However, without appropriate processing tools, this wealth of information remains effectively locked away, unable to contribute to secondary uses such as clinical decision support, quality improvement initiatives, or biomedical research~\citep{wu2022neural, yadav2018mining}.

Natural Language Processing (NLP) has emerged as the key technology for unlocking the potential of clinical text~\citep{wang2018clinical, yadav2018mining}. By enabling computational systems to parse, interpret, and extract meaningful information from human language, NLP bridges the gap between the expressiveness of clinical narratives and the structured data requirements of downstream applications. The field has undergone a remarkable evolution, progressing from early rule-based systems relying on handcrafted patterns and medical vocabularies to sophisticated machine learning approaches, and more recently to transformer-based language models that have achieved unprecedented performance on a wide range of clinical NLP tasks~\citep{kalyan2022ammu, wu2022neural}.

Clinical NLP encompasses a diverse array of tasks that collectively enable the transformation of free text into actionable knowledge. Traditional information extraction tasks include Named Entity Recognition (NER), which identifies and categorizes clinically relevant mentions—diseases, medications, procedures, anatomical sites—within running text~\citep{wang2019clinical, nasar2021named}. Relation Extraction determines the semantic connections between identified entities, such as linking a drug to its dosage or a symptom to its anatomical location~\citep{luo2023systematic}. Temporal reasoning situates clinical events within a timeline, enabling the reconstruction of patient histories~\citep{bethard2017semeval}. Assertion detection and negation handling distinguish between conditions that are present, absent, hypothetical, or attributed to family members~\citep{chapman2013extending}. Text classification supports tasks ranging from document triage to phenotype identification~\citep{yim2016natural}. Together, these capabilities form the foundation for applications such as automated coding, pharmacovigilance, and cohort identification for clinical trials~\citep{kreimeyer2017natural, wang2018clinical}.

The methodological landscape of clinical NLP has evolved through several distinct paradigms. Early systems relied on rule-based approaches, leveraging handcrafted patterns, regular expressions, and medical ontologies such as UMLS and SNOMED-CT to identify clinical concepts~\citep{friedman2004natural}. While interpretable and precise for well-defined patterns, these systems required extensive manual engineering and struggled with the variability inherent in clinical language. The introduction of machine learning brought statistical methods---including Conditional Random Fields (CRFs) and Support Vector Machines---that could learn patterns from annotated corpora, improving robustness to linguistic variation~\citep{jiang2011systematic}. The deep learning revolution, beginning around 2015, saw recurrent neural networks and LSTM-CRF architectures achieve state-of-the-art performance on sequence labeling tasks by learning distributed representations of clinical text~\citep{wu2018deep, si2019enhancing}. However, the most transformative shift came with the introduction of transformer-based pre-trained language models. BERT~\citep{devlin2019bert} and its biomedical and clinical variants---BioBERT~\citep{lee2020biobert}, ClinicalBERT~\citep{huang2019clinicalbert}, and PubMedBERT~\citep{gu2021domain}---demonstrated that self-supervised pre-training on large corpora yields rich contextual representations that transfer effectively to downstream tasks with minimal task-specific data. Scaling these approaches, GatorTron trained on over 90 billion words of text established that larger models and more diverse clinical training data yield consistent improvements across benchmarks~\citep{yang2022large}.

The emergence of Large Language Models (LLMs) has fundamentally expanded the scope of clinical NLP beyond traditional information extraction toward natural language generation and complex reasoning tasks~\citep{singhal2023large, thirunavukarasu2023large}. Clinical text summarization has become a major focus, with LLMs demonstrating the ability to condense lengthy patient records, radiology reports, and doctor-patient dialogues into concise, actionable summaries. Recent studies have shown that adapted LLMs can produce summaries that physicians rate as equivalent or superior to those written by medical experts in terms of completeness, correctness, and conciseness~\citep{vanveen2024adapted}. The automatic generation of discharge summaries represents a particularly impactful application, with several systems demonstrating significant reductions in documentation time while maintaining clinical accuracy. These advances directly address the documentation burden that contributes to physician burnout.

Beyond summarization, LLMs enable entirely new paradigms for clinical decision support. Medical question answering systems can now engage in multi-turn dialogues, providing contextually appropriate responses to complex clinical queries~\citep{singhal2023large}. Diagnostic reasoning represents another frontier, with studies demonstrating that LLMs can match or exceed physician performance on clinical reasoning assessments, though they also exhibit higher rates of incorrect reasoning, underscoring the need for human oversight~\citep{cabral2024clinical}. LLMs have also shown promise for patient-facing applications, including the translation of clinical documentation into patient-friendly language, medication counseling, and preliminary symptom assessment. Crucially, the zero-shot and few-shot learning capabilities of modern LLMs---first demonstrated at scale by GPT-3~\citep{brown2020language}---reduce dependence on large annotated datasets, potentially democratizing access to clinical NLP tools across institutions and languages that lack extensive training resources~\citep{hu2024improving}.

The applications of clinical NLP extend across the entire healthcare continuum. In hospital settings, NLP-powered systems can process millions of clinical documents to extract structured information supporting epidemiological surveillance, quality measurement, and operational analytics~\citep{bean2023hospital}. For clinical research, automated phenotyping algorithms combining NLP with rule-based logic enable the identification of patient cohorts at scales impossible through manual chart review~\citep{kirby2016phekb}. In pharmaceutical development, text mining of both clinical records and biomedical literature accelerates drug discovery, pharmacovigilance, and the identification of drug-drug interactions~\citep{harpaz2014text}. Clinical decision support systems increasingly incorporate NLP components to provide real-time alerts, suggest diagnoses, or recommend evidence-based interventions based on the information contained in clinical notes~\citep{sutton2020overview}.

persistent challenges. The linguistic characteristics of clinical text—including pervasive abbreviations, telegraphic style, misspellings, and domain-specific jargon—differ substantially from the general-domain text on which most NLP models are trained~\citep{dalianis2018clinical}. The sensitive nature of health information imposes strict requirements for data governance, limiting access to training data and constraining the development of robust, generalizable models. Variability across institutions in documentation practices, EHR systems, and clinical workflows means that models trained in one setting may not transfer effectively to another. Furthermore, the high-stakes nature of healthcare demands levels of accuracy, interpretability, and reliability that current systems struggle to consistently achieve~\citep{wu2022neural}. These challenges motivate continued research into methods that can deliver the promise of clinical NLP while respecting the unique constraints of the healthcare domain.

\section{Problem Statement: Data and Privacy, a Fundamental Tension}

% - The paradox: need for massive data to train performant models
% - Regulatory (GDPR, HIPAA) and ethical constraints
% - Limitations of classical approaches (anonymization, restricted access)
% - The hypothesis of synthetic data as a resolution path
The remarkable progress of clinical NLP systems has been driven by a seemingly simple principle: more data yields better models. Deep learning architectures, and particularly large language models, are fundamentally data-hungry, with performance improvements following predictable scaling laws that link model capability to training corpus size. GatorTron's success with 90 billion words of clinical text exemplifies this paradigm---larger models trained on more diverse clinical data consistently outperform their smaller counterparts~\citep{yang2022large}. More recently, Med-PaLM 2 achieved 86.5\% accuracy on the MedQA benchmark through a combination of scale (building on the 540-billion parameter PaLM model) and medical domain fine-tuning, with physician evaluators preferring its answers to those of human doctors on eight of nine clinical axes~\citep{singhal2023medpalm2}. These advances underscore a consistent finding: performance gains in medical AI correlate strongly with access to large-scale, high-quality clinical data. Yet in the medical domain, this appetite for data collides directly with the imperative to protect patient privacy, creating a fundamental tension that shapes and constrains the development of clinical NLP systems.

Health information occupies a privileged position in data protection frameworks worldwide, recognized as among the most sensitive categories of personal data. The European Union's General Data Protection Regulation (GDPR) classifies health data as a ``special category'' under Article 9, subject to a general prohibition on processing with only narrow exceptions for purposes such as medical diagnosis, healthcare provision, or public health research---and even these require explicit legal bases and appropriate safeguards~\citep{gdpr2016regulation}. In the United States, the Health Insurance Portability and Accountability Act (HIPAA) Privacy Rule establishes national standards protecting individually identifiable health information, restricting its use and disclosure while granting patients rights over their medical records~\citep{hipaa1996}. Similar frameworks exist across jurisdictions: France's Loi Informatique et Libertés, as amended to align with GDPR, maintains additional protections for health data processing~\citep{cnil2018loi}, while the UK's Data Protection Act 2018 preserves GDPR-equivalent safeguards post-Brexit with specific provisions for health research~\citep{ukdpa2018}. Sector-specific regulations in healthcare systems globally reinforce these protections. These frameworks share a common recognition that health data merits heightened protection due to the potential for discrimination, stigmatization, and personal harm that could follow from unauthorized disclosure.

Traditional approaches to enabling data sharing while protecting privacy have centered on de-identification: the removal or transformation of identifying information to render records no longer attributable to specific individuals. HIPAA defines two pathways for de-identification: the Safe Harbor method, which requires removal of 18 specific identifier types including names, geographic subdivisions smaller than a state, and dates more specific than year; and the Expert Determination method, which requires a qualified statistician to certify that the risk of re-identification is ``very small''~\citep{hipaa1996}. Similar principles underpin GDPR's concept of anonymization, under which properly anonymized data falls outside the regulation's scope entirely.

However, the reliability of de-identification has been called into question by research demonstrating successful re-identification attacks. A seminal study showed that 87\% of the U.S. population could be uniquely identified by the combination of gender, date of birth, and five-digit ZIP code---information routinely retained in ``de-identified'' datasets~\citep{sweeney2002}. Subsequent work has extended these findings, demonstrating that even datasets de-identified to the HIPAA Safe Harbor standard remain vulnerable when combined with external information sources such as news reports or public records~\citep{elEmam2011}. The re-identification risk is particularly acute for clinical text, where removing explicit identifiers does not eliminate the narrative details---rare conditions, unusual circumstances, distinctive timelines---that can render individuals uniquely identifiable within their communities.

A second limitation of de-identification concerns the trade-off between privacy protection and data utility. Aggressive de-identification that removes or generalizes potentially identifying information inevitably degrades the clinical and linguistic richness that makes the data valuable for NLP research. Dates become years, locations become regions, and specific clinical details may be suppressed entirely. For tasks that depend on temporal reasoning, geographic variation in care patterns, or the precise semantics of clinical narratives, heavily de-identified data may no longer support the intended analyses. This creates a paradox: the more thoroughly data is protected, the less useful it becomes for training and evaluating NLP systems that must ultimately operate on fully detailed clinical text.

The implications for clinical NLP research are profound. Unlike domains where training data can be scraped from the open web or crowdsourced from willing participants, clinical text resides behind institutional firewalls, governed by ethics committees, data use agreements, and legal constraints. Researchers typically cannot access clinical notes from hospitals where they are not employed; even within institutions, data access requires navigating complex governance structures. Cross-institutional collaborations face additional hurdles, as sharing identifiable data across organizational boundaries often requires patient consent or formal data sharing agreements that are slow to negotiate and limited in scope. The result is a fragmented and paradoxical landscape. On one hand, the few publicly available clinical datasets---notably MIMIC and the i2b2/n2c2 challenge corpora---become de facto standards on which most published models are trained and evaluated, despite representing narrow slices of clinical practice (predominantly intensive care units at a single US institution, in the case of MIMIC). Models optimized for these benchmarks may not generalize to the diversity of documentation styles, patient populations, and clinical workflows encountered in real-world settings. On the other hand, institutions that develop models on their own proprietary data often cannot share these models or the datasets underlying them, creating an ecosystem where practical clinical NLP tools remain siloed within individual healthcare systems. This disconnect between academic benchmarks and clinical reality means that published state-of-the-art results may not translate to operational performance, while genuinely useful institutional tools never enter the scientific literature.

These limitations of traditional approaches have motivated growing interest in synthetic data as an alternative paradigm for enabling clinical NLP research while preserving privacy. Synthetic clinical documents are artificially generated texts that mimic the statistical and linguistic properties of real clinical notes without corresponding to any actual patient. When properly generated, synthetic data preserves the semantic richness and domain-specific characteristics that make clinical text valuable for NLP development while severing the link to identifiable individuals that triggers privacy protections.

The potential of synthetic data has been amplified by advances in generative language modeling. Modern LLMs have demonstrated remarkable capability in producing fluent, contextually appropriate text across domains, including medicine. Systems like Synthea pioneered rule-based generation of complete synthetic patient records, while more recent work has leveraged neural language models to generate clinical notes that capture the stylistic and terminological characteristics of authentic documentation~\citep{walonoski2018synthea}. Studies have shown that augmenting real clinical datasets with synthetic examples can improve downstream NLP performance, particularly for minority classes and rare conditions where authentic training data is scarce~\citep{tang2023synthetic, kumichev2024medsyn}. Synthetic data generation may thus address not only privacy constraints but also the class imbalance and data scarcity challenges endemic to clinical NLP.

Yet the promise of synthetic data comes with significant caveats. Unlike general-domain text generation, where plausibility is the primary criterion, clinical synthetic data must satisfy additional requirements of medical coherence, terminological accuracy, and representativeness of actual clinical practice. Generating text that reads like clinical notes is insufficient; the generated content must reflect realistic patient presentations, appropriate clinical reasoning, and accurate relationships between symptoms, diagnoses, and treatments. Furthermore, the very capability that makes LLMs effective generators---their training on vast corpora that may include clinical text---raises questions about whether synthetic outputs might inadvertently leak information about the training data, thereby undermining privacy guarantees.

The quality requirements for clinical NLP training data are particularly demanding. Unlike general-domain NLP, where models can learn from billions of web-scraped tokens of variable quality, clinical NLP benefits disproportionately from high-quality, domain-specific data. The specialized vocabulary, abbreviated writing style, and domain-specific reasoning patterns of clinical text require training corpora that accurately represent these characteristics. Models pre-trained on general text and fine-tuned on clinical data consistently underperform those trained on clinical corpora from the outset~\citep{gu2021domain}. This emphasis on data quality over mere quantity places additional demands on synthetic data generation: synthetic clinical text must not only be private and plentiful but also linguistically and medically authentic.
\section{Research Questions}

% 0.5-1 page
% Formulate 2-3 clear research questions, for example:
% - RQ1: How can we generate realistic and useful synthetic medical documents?
% - RQ2: How can we evaluate the quality and privacy of these documents?
% - RQ3: Can synthetic data effectively train performant downstream models?

\section{Contributions}

% 0.5-1 page
% List of main contributions:
% - A synthetic document generation method based on reward-based generation
% - A multi-criteria evaluation framework (utility, coherence, privacy)
% - A re-identification risk analysis
% - A knowledge distillation method for multilingual clinical entity extraction

\section{Thesis Outline}

% 0.5 page
% Chapter-by-chapter description:
% - Chapter 1: Related Work...
% - Chapter 2: Synthetic Document Generation...
% - Chapter 3: Improvements and Privacy-Preserving Evaluation...
% - Chapter 4: Knowledge Distillation for Clinical Entity Extraction...
